%%%%%%%%%%%%%%%%%%%%%%%%%%%%%%%%%
%
%       EXERCÍCIO
%
%%%%%%%%%%%%%%%%%%%%%%%%%%%%%%%%%

\ifdefstring{\atividade}{prova}{%
    \ifdefstring{\modo}{objetivo}{%
        \renewcommand{\valorquestao}{\ValorQObj\ ponto}
    }{%
        \renewcommand{\valorquestao}{\ValorQDisc\ pontos}
    }%
}%

\begin{exercicioBanco}[\valorquestao]
Assinale verdadeira (V) ou falso (F) a cada uma das seguintes sentenças.
%
\begin{enumerate}[\(\bullet\)]
    \item A função \(f(x) = 5 - 3x\) é crescente.
    %
    \item O zero de \(f(x) = -x + 4\) é igual a \(4\).
    %
    \item A função \(f(x) = x + 2\) é negativa no intervalo \(]-\infty,-2[\).
    %
    \item O coeficiente angular de \(f(x) = 2x - 1\) é igual a \(-1\).
\end{enumerate}

% Define as alternativas
\newcommand{\alternativas}{%
    \begin{center}
        \begin{tabularx}{\textwidth}{XXXXX}
            (a) (F, V, V, V). &
            (b) (F, V, F, F). &
            (c) (V, V, V, F). &
            (d) (F, F, V, F). &
            (e) (F, V, V, F).
        \end{tabularx}
    \end{center}
}

% Define a resposta correta
\newcommand{\resposta}{E}

% Lógica condicional para exibição
\ifdefstring{\atividade}{lista}{%
    \alternativas
    \vspace{0.5em}
    
    \noindent\textbf{Resposta:} letra \textbf{\resposta}.
}{%
    \ifdefstring{\modo}{objetivo}{%
        \alternativas
    }{%
        % modo = discursiva → não mostra alternativas
    }
}
\end{exercicioBanco}

